% ============================================================
% DROP-IN SUBSECTION FOR report.tex
%
% Placement:
%   Insert at the end of Section 5 "Exploratory Autoencoder Development",
%   immediately after the subsection labelled sec:ali_baseline and before
%   Section 6 "The Final Keras CAE (Milestone 5)".
%
% Recommended include from report.tex:
%   % ============================================================
% DROP-IN SUBSECTION FOR report.tex
%
% Placement:
%   Insert at the end of Section 5 "Exploratory Autoencoder Development",
%   immediately after the subsection labelled sec:ali_baseline and before
%   Section 6 "The Final Keras CAE (Milestone 5)".
%
% Recommended include from report.tex:
%   % ============================================================
% DROP-IN SUBSECTION FOR report.tex
%
% Placement:
%   Insert at the end of Section 5 "Exploratory Autoencoder Development",
%   immediately after the subsection labelled sec:ali_baseline and before
%   Section 6 "The Final Keras CAE (Milestone 5)".
%
% Recommended include from report.tex:
%   % ============================================================
% DROP-IN SUBSECTION FOR report.tex
%
% Placement:
%   Insert at the end of Section 5 "Exploratory Autoencoder Development",
%   immediately after the subsection labelled sec:ali_baseline and before
%   Section 6 "The Final Keras CAE (Milestone 5)".
%
% Recommended include from report.tex:
%   \input{sections/screw_exploration}
%
% Figures:
%   figures/screw/fig_v2_v4_residual_snr.png
%   figures/screw/fig_v2_localization_examples.png
%   figures/screw/fig_v3_masking_bug_vs_fix.png
%   figures/screw/figures/masking_example_background_miss.png
%   figures/screw/figures/masking_example_object_hit.png
%
% Cross-references below use the existing report labels rather than
% hardcoded section/table numbers.
%
% Note for the main report:
%   Section 2.2 / 4.2 still need the separate orientation wording update
%   so that screw is treated as a random-rotation object-category exception.
% ============================================================

\subsection{Parallel Screw-Category Exploration}
\label{subsec:screw-exploration}

Alongside the bottle-category exploration described above, we conducted a
parallel iterative development phase on the screw category. Screw presents
two structural challenges not shared by bottle: its images are
single-channel greyscale rather than RGB, and, per the original MVTec AD
benchmark, its parts are photographed at random rotation on the inspection
surface rather than arriving in a consistent orientation, one of the
reasons the benchmark identifies screw as one of the dataset's most
difficult categories. Under the same deterministic 85\%/15\% protocol used
throughout this report, screw's 320 normal training images are divided
into 272 fitting images and 48 validation images. As with the bottle
milestones above, the official test set (160 images: 41 normal, 119
defective) is used to score each milestone as it is built, not reserved
for a single blind check at the end, so the results below are exploratory
and test-informed in exactly the same sense as Milestones 1-4 for bottle.
We report this phase for the same reason we report those milestones: the
dead ends are as informative as the improvements.

This screw exploration was implemented independently in PyTorch, separately
from the Keras implementation of the final CAE described in
Section~\ref{sec:cae}. The two are not architecturally identical and share no code.

\textbf{A note on naming.} The five variants below are referred to as
Milestones S1-S5 in this write-up, to parallel this report's Milestone
1-5 naming for bottle. The figures accompanying Milestones S3-S5 were
generated directly by the exploration notebook and keep that notebook's
own labels (V1, V2, V3, V3b, V4) baked into the image content. The
mapping is: V1 = S1, V2 = S2, V3 = S3, V3b = S4, V4 = S5. 

Supplementary visualisations from these experiments are collected in
Appendix~\ref{sec:screw-exploration-figures}.

\subsubsection{Baselines: Z-Score Model and the Accuracy Paradox in Practice}

Before any learned model, we built a non-learned baseline: a per-pixel mean and
standard deviation computed from the normal training images, with test-time
pixels flagged when they deviate significantly from that per-pixel statistic,
conceptually identical to the Variation Model benchmarked in the original
MVTec AD paper. A dedicated three-seed robustness check (seeds 42, 17, 100)
confirmed this baseline is stable regardless of the train/validation split
seed, with a Pixel AP standard deviation of 0.0003 and a Pixel AUROC standard
deviation of 0.0018; mean Pixel AP across the three seeds is 0.0123 at 98.43\%
accuracy.

To make the accuracy-paradox argument from Section~\ref{sec:imbalance} concrete on this
category rather than only cited, we additionally scored a trivial
``always normal'' classifier against the true defect masks: it achieved
99.75\% pixel accuracy against just 0.0025 AP. The z-score baseline
deliberately trades a little accuracy for a large jump in AP, the opposite
direction accuracy alone would suggest: under the single seed-42 run used
throughout the rest of this section, it reaches 98.43\% accuracy and 0.0126
AP, consistent with the three-seed mean above, which is the clearest
demonstration in our screw experiments of why AP-style metrics are
necessary here.

\subsubsection{Milestone S1: Baseline Convolutional Autoencoder}

Our first screw architecture used a standard encoder-decoder with an implicit
spatial bottleneck (no explicit dense compression layer), trained with MSE
reconstruction loss on normal-only images. This establishes the reference
point every later change is measured against. The result was a Pixel AP of
0.0225, a real but modest signal that confirmed the pipeline worked end to end
without yet being competitive.

\subsubsection{Milestone S2: Explicit 128-Dimensional Bottleneck}

We added a dense linear bottleneck, forcing a much stronger, more deliberate
compression than Milestone S1's implicit spatial one. The result was a Pixel
AP of 0.0453, roughly double Milestone S1's. This was the first strong
confirmation on screw that compression strength is a real, controllable lever
for anomaly separation, not just capacity for cleaner reconstruction, the same
conclusion Section~\ref{sec:ae3} reaches independently for bottle.

\subsubsection{Milestone S3: Masked-Restoration Self-Supervision and the Overlap Diagnostic}

We tested a different strategy for the same underlying goal explored in
Milestone S2, preventing the model from trivially learning to copy its input,
adapted from the masked-restoration technique introduced for bottle in
Milestone 2 (Section~\ref{sec:ae2}): random patches were zeroed and the model
was trained to restore the full image. This failed decisively
on screw. We diagnosed the failure quantitatively rather than assuming it: a
targeted check sampling 500 random $32\times32$ training patches found that
only 190 (38.0\%) touched the screw object at all, meaning 62.0\% of masked
patches during training were pure background. Most of the training signal
taught the model nothing about screw structure. The result was a Pixel AP of
0.0117, the worst model in this phase, below even the non-learned z-score
baseline (0.0123).

Milestone S3 uses Python's built-in random module for masked-patch placement.
It calls \texttt{random.randint()}, which is not controlled by
\texttt{torch.manual\_seed()}. Because \texttt{random.seed()} is never set before this milestone trains, its patch
locations, and therefore its trained weights and final metrics, depend on
whatever random state the kernel happens to start with: across two
independent reruns of identical code, Pixel AP ranged from 0.0146 to 0.0394
and image-level AUROC ranged from 0.000 to 0.563. Milestone S4 (below) uses
the same unseeded function but happens to run after an unrelated
diagnostic cell that calls \texttt{random.seed(0)} as a side effect, which
is why it reproduces exactly across reruns despite sharing the same
underlying gap; that reproducibility is an artefact of cell execution
order, not an intentional seed. The fix is to add
\texttt{random.seed(RANDOM\_SEED)} alongside the existing PyTorch seed.

Because this milestone's patch placement is unseeded, where the mask lands
varies from run to run rather than being fixed by the model or data. Two
additional single-patch examples illustrate this directly: in one run the
masked patch sits almost entirely on background, away from the screw body;
in another run of the identical code, the patch lands closer to the object,
overlapping its edge. This is the same unseeded call landing differently
each time the kernel starts, not two different bugs, which is exactly why
this milestone's metrics vary run to run rather than the placement being
consistently bad or consistently fine.

\subsubsection{Milestone S4: Object-Centred Masking Fix}

We directly addressed the diagnosed Milestone S3 bug by restricting mask
placement to the object's foreground region, using the same
foreground-extraction pipeline described in Section~\ref{sec:foreground}, rather than masking
uniformly at random across the full frame. The bug fix itself is correct and
verifiable independent of any model score. However, fixing the bug did not
translate into a competitive result, but it did help: this model's Pixel AP
(0.0173) is roughly 1.5x Milestone S3's (0.0117) and moves back above the
non-learned z-score baseline (0.0123), though it still falls well short of
the bottleneck-based milestones (S1, S2, S5). We consider this an important
negative-then-partial-recovery result of our screw experiments: it
demonstrates that a correct bug diagnosis and a correct fix do not
automatically make the underlying method competitive.

\subsubsection{Milestone S5: L1 Loss Ablation}

Our final screw variant kept Milestone S2's architecture unchanged and
replaced the MSE reconstruction loss with an L1 (mean absolute error) loss,
motivated by the hypothesis that L1's sharper, less blur-prone gradient
behaviour would produce cleaner separation between normal and defective pixel
errors. This produced the best result of the screw exploration phase: Pixel
AP 0.0581, Pixel AUROC 0.9075, and a top-1\% defect-pixel capture rate of
0.2508.

To rule out the simplest alternative explanation, that L1 was simply
elevating error everywhere including the background rather than genuinely
separating defect regions better, we measured mean reconstruction error
inside and outside the ground-truth defect mask across all 119 defective
screw test images (Table~\ref{tab:screw-snr}). The L1 model's background
error is only marginally higher than the MSE model's, while its
defect-region error is proportionally much higher, giving it a higher
signal-to-noise ratio and ruling out the simple noise-amplification
explanation. We note this analysis was run against an earlier,
pre-final checkpoint pair, so the qualitative finding is suggestive rather
than confirmed against the exact final models in
Table~\ref{tab:screw-results}.

\begin{table}[h]
\centering
\caption{Mean reconstruction error inside and outside the ground-truth defect
mask, Milestone S2 (MSE) vs.\ Milestone S5 (L1), across all 119 defective
screw test images. Checkpoints predate the final seed-42 models in
Table~\ref{tab:screw-results}.}
\begin{tabular}{lccc}
\toprule
Model & Defect Error (signal) & Background Error (noise) & SNR \\
\midrule
Milestone S2 (MSE) & 0.018015 & 0.001330 & 13.54 \\
Milestone S5 (L1 loss) & 0.024089 & 0.001447 & 16.65 \\
\bottomrule
\end{tabular}
\label{tab:screw-snr}
\end{table}

\subsubsection{Summary of Screw Exploration Results}

Table~\ref{tab:screw-results} summarises the full screw exploration phase.
All autoencoder rows use seed 42, a 20-epoch training budget, and the 85/15
split described above; the baseline row is the mean of the three-seed check.
Pixel F1 and Image F1 use the same validation-derived thresholds as the
adjacent accuracy columns (99th percentile pixel-level, 95th percentile
image-level). Milestones S3 and S4's values should be read as a single-run
snapshot affected by the diagnosed seed-instability issue in their masking
implementation (unseeded \texttt{random.randint()}, not covered by
\texttt{torch.manual\_seed()}): as Milestone S3's write-up above shows, a
different unseeded rerun of the same code can produce materially different
numbers for this specific pair. Every other row in this table is seed-stable.

\begin{table}[h]
\centering
\caption{Screw-category results across the exploratory autoencoder variants.}
\resizebox{\textwidth}{!}{%
\begin{tabular}{lcccccccc}
\toprule
Model & Pixel AP & Pixel Acc & Pixel F1 & Top-1\% Capture & Pixel AUROC & Img F1 & Img MSE AUROC & Top-K AUROC \\
\midrule
Z-score baseline & 0.0123 & 0.9843 & -- & -- & 0.6887 & -- & -- & -- \\
Milestone S1 & 0.0225 & 0.9885 & 0.0557 & 0.1417 & 0.8723 & 0.000 & 0.8541 & 0.9561 \\
Milestone S2 & 0.0453 & 0.9872 & 0.0900 & 0.2269 & 0.8849 & 0.168 & 0.5431 & 0.7225 \\
Milestone S3 & 0.0117 & 0.9868 & 0.0320 & 0.0852 & 0.7765 & 0.025 & 0.0047 & 0.3210 \\
Milestone S4 & 0.0173 & 0.9826 & 0.0526 & 0.2015 & 0.7401 & 0.053 & 0.0842 & 0.0502 \\
\textbf{Milestone S5 (L1 loss)} & \textbf{0.0581} & 0.9871 & 0.1029 & \textbf{0.2508} & \textbf{0.9075} & 0.126 & 0.8176 & 0.8721 \\
\bottomrule
\end{tabular}%
}
\label{tab:screw-results}
\end{table}

Image F1 of 0.000 for Milestone S1 is a genuine measurement at the
validation-derived operating threshold, not a data error. It means that this
operating point produced no true-positive image detections on the screw test
set. The result should therefore be interpreted together with the
threshold-independent image-level AUROC, which shows that the model can retain
ranking information even when its chosen operating threshold performs poorly.

These exploratory results predate and motivated the production PatchCore and
Keras CAE baseline comparison reported in Section~\ref{sec:evaluation}, where screw achieves a
PatchCore image F1 of 0.796 and pixel AUPIMO of 0.380, and a Keras CAE image
F1 of 0.421 and pixel AUPIMO of 0.005 (Table~\ref{tab:final_results}), figures that already reflect
the lessons of this exploration: screw's rotational variability and greyscale
presentation make it one of the harder categories in the benchmark,
consistent with what we found here.

The gap is largest on image-level F1 specifically: Milestone S5, the best
model in this exploration, reaches an image F1 of only 0.126, well short of
both the PatchCore (0.796) and Keras CAE (0.421) screw baselines above. This
is expected rather than a weakness specific to screw: none of these
exploratory autoencoders use a pretrained backbone, and the same
from-scratch-versus-pretrained-features gap is already visible in this
report's own bottle exploration (Section~\ref{sec:ali_baseline}), where a frozen-ImageNet
5-NN baseline (AUROC 0.988) decisively outperformed every from-scratch
autoencoder variant (best AUROC 0.877) and directly motivated adopting
PatchCore. The screw milestones here should be read the same way: as
exploratory evidence about masking, compression, and loss-function choices
on a from-scratch architecture, not as a competitive alternative to the
production models above.

%
% Figures:
%   figures/screw/fig_v2_v4_residual_snr.png
%   figures/screw/fig_v2_localization_examples.png
%   figures/screw/fig_v3_masking_bug_vs_fix.png
%   figures/screw/figures/masking_example_background_miss.png
%   figures/screw/figures/masking_example_object_hit.png
%
% Cross-references below use the existing report labels rather than
% hardcoded section/table numbers.
%
% Note for the main report:
%   Section 2.2 / 4.2 still need the separate orientation wording update
%   so that screw is treated as a random-rotation object-category exception.
% ============================================================

\subsection{Parallel Screw-Category Exploration}
\label{subsec:screw-exploration}

Alongside the bottle-category exploration described above, we conducted a
parallel iterative development phase on the screw category. Screw presents
two structural challenges not shared by bottle: its images are
single-channel greyscale rather than RGB, and, per the original MVTec AD
benchmark, its parts are photographed at random rotation on the inspection
surface rather than arriving in a consistent orientation, one of the
reasons the benchmark identifies screw as one of the dataset's most
difficult categories. Under the same deterministic 85\%/15\% protocol used
throughout this report, screw's 320 normal training images are divided
into 272 fitting images and 48 validation images. As with the bottle
milestones above, the official test set (160 images: 41 normal, 119
defective) is used to score each milestone as it is built, not reserved
for a single blind check at the end, so the results below are exploratory
and test-informed in exactly the same sense as Milestones 1-4 for bottle.
We report this phase for the same reason we report those milestones: the
dead ends are as informative as the improvements.

This screw exploration was implemented independently in PyTorch, separately
from the Keras implementation of the final CAE described in
Section~\ref{sec:cae}. The two are not architecturally identical and share no code.

\textbf{A note on naming.} The five variants below are referred to as
Milestones S1-S5 in this write-up, to parallel this report's Milestone
1-5 naming for bottle. The figures accompanying Milestones S3-S5 were
generated directly by the exploration notebook and keep that notebook's
own labels (V1, V2, V3, V3b, V4) baked into the image content. The
mapping is: V1 = S1, V2 = S2, V3 = S3, V3b = S4, V4 = S5. 

Supplementary visualisations from these experiments are collected in
Appendix~\ref{sec:screw-exploration-figures}.

\subsubsection{Baselines: Z-Score Model and the Accuracy Paradox in Practice}

Before any learned model, we built a non-learned baseline: a per-pixel mean and
standard deviation computed from the normal training images, with test-time
pixels flagged when they deviate significantly from that per-pixel statistic,
conceptually identical to the Variation Model benchmarked in the original
MVTec AD paper. A dedicated three-seed robustness check (seeds 42, 17, 100)
confirmed this baseline is stable regardless of the train/validation split
seed, with a Pixel AP standard deviation of 0.0003 and a Pixel AUROC standard
deviation of 0.0018; mean Pixel AP across the three seeds is 0.0123 at 98.43\%
accuracy.

To make the accuracy-paradox argument from Section~\ref{sec:imbalance} concrete on this
category rather than only cited, we additionally scored a trivial
``always normal'' classifier against the true defect masks: it achieved
99.75\% pixel accuracy against just 0.0025 AP. The z-score baseline
deliberately trades a little accuracy for a large jump in AP, the opposite
direction accuracy alone would suggest: under the single seed-42 run used
throughout the rest of this section, it reaches 98.43\% accuracy and 0.0126
AP, consistent with the three-seed mean above, which is the clearest
demonstration in our screw experiments of why AP-style metrics are
necessary here.

\subsubsection{Milestone S1: Baseline Convolutional Autoencoder}

Our first screw architecture used a standard encoder-decoder with an implicit
spatial bottleneck (no explicit dense compression layer), trained with MSE
reconstruction loss on normal-only images. This establishes the reference
point every later change is measured against. The result was a Pixel AP of
0.0225, a real but modest signal that confirmed the pipeline worked end to end
without yet being competitive.

\subsubsection{Milestone S2: Explicit 128-Dimensional Bottleneck}

We added a dense linear bottleneck, forcing a much stronger, more deliberate
compression than Milestone S1's implicit spatial one. The result was a Pixel
AP of 0.0453, roughly double Milestone S1's. This was the first strong
confirmation on screw that compression strength is a real, controllable lever
for anomaly separation, not just capacity for cleaner reconstruction, the same
conclusion Section~\ref{sec:ae3} reaches independently for bottle.

\subsubsection{Milestone S3: Masked-Restoration Self-Supervision and the Overlap Diagnostic}

We tested a different strategy for the same underlying goal explored in
Milestone S2, preventing the model from trivially learning to copy its input,
adapted from the masked-restoration technique introduced for bottle in
Milestone 2 (Section~\ref{sec:ae2}): random patches were zeroed and the model
was trained to restore the full image. This failed decisively
on screw. We diagnosed the failure quantitatively rather than assuming it: a
targeted check sampling 500 random $32\times32$ training patches found that
only 190 (38.0\%) touched the screw object at all, meaning 62.0\% of masked
patches during training were pure background. Most of the training signal
taught the model nothing about screw structure. The result was a Pixel AP of
0.0117, the worst model in this phase, below even the non-learned z-score
baseline (0.0123).

Milestone S3 uses Python's built-in random module for masked-patch placement.
It calls \texttt{random.randint()}, which is not controlled by
\texttt{torch.manual\_seed()}. Because \texttt{random.seed()} is never set before this milestone trains, its patch
locations, and therefore its trained weights and final metrics, depend on
whatever random state the kernel happens to start with: across two
independent reruns of identical code, Pixel AP ranged from 0.0146 to 0.0394
and image-level AUROC ranged from 0.000 to 0.563. Milestone S4 (below) uses
the same unseeded function but happens to run after an unrelated
diagnostic cell that calls \texttt{random.seed(0)} as a side effect, which
is why it reproduces exactly across reruns despite sharing the same
underlying gap; that reproducibility is an artefact of cell execution
order, not an intentional seed. The fix is to add
\texttt{random.seed(RANDOM\_SEED)} alongside the existing PyTorch seed.

Because this milestone's patch placement is unseeded, where the mask lands
varies from run to run rather than being fixed by the model or data. Two
additional single-patch examples illustrate this directly: in one run the
masked patch sits almost entirely on background, away from the screw body;
in another run of the identical code, the patch lands closer to the object,
overlapping its edge. This is the same unseeded call landing differently
each time the kernel starts, not two different bugs, which is exactly why
this milestone's metrics vary run to run rather than the placement being
consistently bad or consistently fine.

\subsubsection{Milestone S4: Object-Centred Masking Fix}

We directly addressed the diagnosed Milestone S3 bug by restricting mask
placement to the object's foreground region, using the same
foreground-extraction pipeline described in Section~\ref{sec:foreground}, rather than masking
uniformly at random across the full frame. The bug fix itself is correct and
verifiable independent of any model score. However, fixing the bug did not
translate into a competitive result, but it did help: this model's Pixel AP
(0.0173) is roughly 1.5x Milestone S3's (0.0117) and moves back above the
non-learned z-score baseline (0.0123), though it still falls well short of
the bottleneck-based milestones (S1, S2, S5). We consider this an important
negative-then-partial-recovery result of our screw experiments: it
demonstrates that a correct bug diagnosis and a correct fix do not
automatically make the underlying method competitive.

\subsubsection{Milestone S5: L1 Loss Ablation}

Our final screw variant kept Milestone S2's architecture unchanged and
replaced the MSE reconstruction loss with an L1 (mean absolute error) loss,
motivated by the hypothesis that L1's sharper, less blur-prone gradient
behaviour would produce cleaner separation between normal and defective pixel
errors. This produced the best result of the screw exploration phase: Pixel
AP 0.0581, Pixel AUROC 0.9075, and a top-1\% defect-pixel capture rate of
0.2508.

To rule out the simplest alternative explanation, that L1 was simply
elevating error everywhere including the background rather than genuinely
separating defect regions better, we measured mean reconstruction error
inside and outside the ground-truth defect mask across all 119 defective
screw test images (Table~\ref{tab:screw-snr}). The L1 model's background
error is only marginally higher than the MSE model's, while its
defect-region error is proportionally much higher, giving it a higher
signal-to-noise ratio and ruling out the simple noise-amplification
explanation. We note this analysis was run against an earlier,
pre-final checkpoint pair, so the qualitative finding is suggestive rather
than confirmed against the exact final models in
Table~\ref{tab:screw-results}.

\begin{table}[h]
\centering
\caption{Mean reconstruction error inside and outside the ground-truth defect
mask, Milestone S2 (MSE) vs.\ Milestone S5 (L1), across all 119 defective
screw test images. Checkpoints predate the final seed-42 models in
Table~\ref{tab:screw-results}.}
\begin{tabular}{lccc}
\toprule
Model & Defect Error (signal) & Background Error (noise) & SNR \\
\midrule
Milestone S2 (MSE) & 0.018015 & 0.001330 & 13.54 \\
Milestone S5 (L1 loss) & 0.024089 & 0.001447 & 16.65 \\
\bottomrule
\end{tabular}
\label{tab:screw-snr}
\end{table}

\subsubsection{Summary of Screw Exploration Results}

Table~\ref{tab:screw-results} summarises the full screw exploration phase.
All autoencoder rows use seed 42, a 20-epoch training budget, and the 85/15
split described above; the baseline row is the mean of the three-seed check.
Pixel F1 and Image F1 use the same validation-derived thresholds as the
adjacent accuracy columns (99th percentile pixel-level, 95th percentile
image-level). Milestones S3 and S4's values should be read as a single-run
snapshot affected by the diagnosed seed-instability issue in their masking
implementation (unseeded \texttt{random.randint()}, not covered by
\texttt{torch.manual\_seed()}): as Milestone S3's write-up above shows, a
different unseeded rerun of the same code can produce materially different
numbers for this specific pair. Every other row in this table is seed-stable.

\begin{table}[h]
\centering
\caption{Screw-category results across the exploratory autoencoder variants.}
\resizebox{\textwidth}{!}{%
\begin{tabular}{lcccccccc}
\toprule
Model & Pixel AP & Pixel Acc & Pixel F1 & Top-1\% Capture & Pixel AUROC & Img F1 & Img MSE AUROC & Top-K AUROC \\
\midrule
Z-score baseline & 0.0123 & 0.9843 & -- & -- & 0.6887 & -- & -- & -- \\
Milestone S1 & 0.0225 & 0.9885 & 0.0557 & 0.1417 & 0.8723 & 0.000 & 0.8541 & 0.9561 \\
Milestone S2 & 0.0453 & 0.9872 & 0.0900 & 0.2269 & 0.8849 & 0.168 & 0.5431 & 0.7225 \\
Milestone S3 & 0.0117 & 0.9868 & 0.0320 & 0.0852 & 0.7765 & 0.025 & 0.0047 & 0.3210 \\
Milestone S4 & 0.0173 & 0.9826 & 0.0526 & 0.2015 & 0.7401 & 0.053 & 0.0842 & 0.0502 \\
\textbf{Milestone S5 (L1 loss)} & \textbf{0.0581} & 0.9871 & 0.1029 & \textbf{0.2508} & \textbf{0.9075} & 0.126 & 0.8176 & 0.8721 \\
\bottomrule
\end{tabular}%
}
\label{tab:screw-results}
\end{table}

Image F1 of 0.000 for Milestone S1 is a genuine measurement at the
validation-derived operating threshold, not a data error. It means that this
operating point produced no true-positive image detections on the screw test
set. The result should therefore be interpreted together with the
threshold-independent image-level AUROC, which shows that the model can retain
ranking information even when its chosen operating threshold performs poorly.

These exploratory results predate and motivated the production PatchCore and
Keras CAE baseline comparison reported in Section~\ref{sec:evaluation}, where screw achieves a
PatchCore image F1 of 0.796 and pixel AUPIMO of 0.380, and a Keras CAE image
F1 of 0.421 and pixel AUPIMO of 0.005 (Table~\ref{tab:final_results}), figures that already reflect
the lessons of this exploration: screw's rotational variability and greyscale
presentation make it one of the harder categories in the benchmark,
consistent with what we found here.

The gap is largest on image-level F1 specifically: Milestone S5, the best
model in this exploration, reaches an image F1 of only 0.126, well short of
both the PatchCore (0.796) and Keras CAE (0.421) screw baselines above. This
is expected rather than a weakness specific to screw: none of these
exploratory autoencoders use a pretrained backbone, and the same
from-scratch-versus-pretrained-features gap is already visible in this
report's own bottle exploration (Section~\ref{sec:ali_baseline}), where a frozen-ImageNet
5-NN baseline (AUROC 0.988) decisively outperformed every from-scratch
autoencoder variant (best AUROC 0.877) and directly motivated adopting
PatchCore. The screw milestones here should be read the same way: as
exploratory evidence about masking, compression, and loss-function choices
on a from-scratch architecture, not as a competitive alternative to the
production models above.

%
% Figures:
%   figures/screw/fig_v2_v4_residual_snr.png
%   figures/screw/fig_v2_localization_examples.png
%   figures/screw/fig_v3_masking_bug_vs_fix.png
%   figures/screw/figures/masking_example_background_miss.png
%   figures/screw/figures/masking_example_object_hit.png
%
% Cross-references below use the existing report labels rather than
% hardcoded section/table numbers.
%
% Note for the main report:
%   Section 2.2 / 4.2 still need the separate orientation wording update
%   so that screw is treated as a random-rotation object-category exception.
% ============================================================

\subsection{Parallel Screw-Category Exploration}
\label{subsec:screw-exploration}

Alongside the bottle-category exploration described above, we conducted a
parallel iterative development phase on the screw category. Screw presents
two structural challenges not shared by bottle: its images are
single-channel greyscale rather than RGB, and, per the original MVTec AD
benchmark, its parts are photographed at random rotation on the inspection
surface rather than arriving in a consistent orientation, one of the
reasons the benchmark identifies screw as one of the dataset's most
difficult categories. Under the same deterministic 85\%/15\% protocol used
throughout this report, screw's 320 normal training images are divided
into 272 fitting images and 48 validation images. As with the bottle
milestones above, the official test set (160 images: 41 normal, 119
defective) is used to score each milestone as it is built, not reserved
for a single blind check at the end, so the results below are exploratory
and test-informed in exactly the same sense as Milestones 1-4 for bottle.
We report this phase for the same reason we report those milestones: the
dead ends are as informative as the improvements.

This screw exploration was implemented independently in PyTorch, separately
from the Keras implementation of the final CAE described in
Section~\ref{sec:cae}. The two are not architecturally identical and share no code.

\textbf{A note on naming.} The five variants below are referred to as
Milestones S1-S5 in this write-up, to parallel this report's Milestone
1-5 naming for bottle. The figures accompanying Milestones S3-S5 were
generated directly by the exploration notebook and keep that notebook's
own labels (V1, V2, V3, V3b, V4) baked into the image content. The
mapping is: V1 = S1, V2 = S2, V3 = S3, V3b = S4, V4 = S5. 

Supplementary visualisations from these experiments are collected in
Appendix~\ref{sec:screw-exploration-figures}.

\subsubsection{Baselines: Z-Score Model and the Accuracy Paradox in Practice}

Before any learned model, we built a non-learned baseline: a per-pixel mean and
standard deviation computed from the normal training images, with test-time
pixels flagged when they deviate significantly from that per-pixel statistic,
conceptually identical to the Variation Model benchmarked in the original
MVTec AD paper. A dedicated three-seed robustness check (seeds 42, 17, 100)
confirmed this baseline is stable regardless of the train/validation split
seed, with a Pixel AP standard deviation of 0.0003 and a Pixel AUROC standard
deviation of 0.0018; mean Pixel AP across the three seeds is 0.0123 at 98.43\%
accuracy.

To make the accuracy-paradox argument from Section~\ref{sec:imbalance} concrete on this
category rather than only cited, we additionally scored a trivial
``always normal'' classifier against the true defect masks: it achieved
99.75\% pixel accuracy against just 0.0025 AP. The z-score baseline
deliberately trades a little accuracy for a large jump in AP, the opposite
direction accuracy alone would suggest: under the single seed-42 run used
throughout the rest of this section, it reaches 98.43\% accuracy and 0.0126
AP, consistent with the three-seed mean above, which is the clearest
demonstration in our screw experiments of why AP-style metrics are
necessary here.

\subsubsection{Milestone S1: Baseline Convolutional Autoencoder}

Our first screw architecture used a standard encoder-decoder with an implicit
spatial bottleneck (no explicit dense compression layer), trained with MSE
reconstruction loss on normal-only images. This establishes the reference
point every later change is measured against. The result was a Pixel AP of
0.0225, a real but modest signal that confirmed the pipeline worked end to end
without yet being competitive.

\subsubsection{Milestone S2: Explicit 128-Dimensional Bottleneck}

We added a dense linear bottleneck, forcing a much stronger, more deliberate
compression than Milestone S1's implicit spatial one. The result was a Pixel
AP of 0.0453, roughly double Milestone S1's. This was the first strong
confirmation on screw that compression strength is a real, controllable lever
for anomaly separation, not just capacity for cleaner reconstruction, the same
conclusion Section~\ref{sec:ae3} reaches independently for bottle.

\subsubsection{Milestone S3: Masked-Restoration Self-Supervision and the Overlap Diagnostic}

We tested a different strategy for the same underlying goal explored in
Milestone S2, preventing the model from trivially learning to copy its input,
adapted from the masked-restoration technique introduced for bottle in
Milestone 2 (Section~\ref{sec:ae2}): random patches were zeroed and the model
was trained to restore the full image. This failed decisively
on screw. We diagnosed the failure quantitatively rather than assuming it: a
targeted check sampling 500 random $32\times32$ training patches found that
only 190 (38.0\%) touched the screw object at all, meaning 62.0\% of masked
patches during training were pure background. Most of the training signal
taught the model nothing about screw structure. The result was a Pixel AP of
0.0117, the worst model in this phase, below even the non-learned z-score
baseline (0.0123).

Milestone S3 uses Python's built-in random module for masked-patch placement.
It calls \texttt{random.randint()}, which is not controlled by
\texttt{torch.manual\_seed()}. Because \texttt{random.seed()} is never set before this milestone trains, its patch
locations, and therefore its trained weights and final metrics, depend on
whatever random state the kernel happens to start with: across two
independent reruns of identical code, Pixel AP ranged from 0.0146 to 0.0394
and image-level AUROC ranged from 0.000 to 0.563. Milestone S4 (below) uses
the same unseeded function but happens to run after an unrelated
diagnostic cell that calls \texttt{random.seed(0)} as a side effect, which
is why it reproduces exactly across reruns despite sharing the same
underlying gap; that reproducibility is an artefact of cell execution
order, not an intentional seed. The fix is to add
\texttt{random.seed(RANDOM\_SEED)} alongside the existing PyTorch seed.

Because this milestone's patch placement is unseeded, where the mask lands
varies from run to run rather than being fixed by the model or data. Two
additional single-patch examples illustrate this directly: in one run the
masked patch sits almost entirely on background, away from the screw body;
in another run of the identical code, the patch lands closer to the object,
overlapping its edge. This is the same unseeded call landing differently
each time the kernel starts, not two different bugs, which is exactly why
this milestone's metrics vary run to run rather than the placement being
consistently bad or consistently fine.

\subsubsection{Milestone S4: Object-Centred Masking Fix}

We directly addressed the diagnosed Milestone S3 bug by restricting mask
placement to the object's foreground region, using the same
foreground-extraction pipeline described in Section~\ref{sec:foreground}, rather than masking
uniformly at random across the full frame. The bug fix itself is correct and
verifiable independent of any model score. However, fixing the bug did not
translate into a competitive result, but it did help: this model's Pixel AP
(0.0173) is roughly 1.5x Milestone S3's (0.0117) and moves back above the
non-learned z-score baseline (0.0123), though it still falls well short of
the bottleneck-based milestones (S1, S2, S5). We consider this an important
negative-then-partial-recovery result of our screw experiments: it
demonstrates that a correct bug diagnosis and a correct fix do not
automatically make the underlying method competitive.

\subsubsection{Milestone S5: L1 Loss Ablation}

Our final screw variant kept Milestone S2's architecture unchanged and
replaced the MSE reconstruction loss with an L1 (mean absolute error) loss,
motivated by the hypothesis that L1's sharper, less blur-prone gradient
behaviour would produce cleaner separation between normal and defective pixel
errors. This produced the best result of the screw exploration phase: Pixel
AP 0.0581, Pixel AUROC 0.9075, and a top-1\% defect-pixel capture rate of
0.2508.

To rule out the simplest alternative explanation, that L1 was simply
elevating error everywhere including the background rather than genuinely
separating defect regions better, we measured mean reconstruction error
inside and outside the ground-truth defect mask across all 119 defective
screw test images (Table~\ref{tab:screw-snr}). The L1 model's background
error is only marginally higher than the MSE model's, while its
defect-region error is proportionally much higher, giving it a higher
signal-to-noise ratio and ruling out the simple noise-amplification
explanation. We note this analysis was run against an earlier,
pre-final checkpoint pair, so the qualitative finding is suggestive rather
than confirmed against the exact final models in
Table~\ref{tab:screw-results}.

\begin{table}[h]
\centering
\caption{Mean reconstruction error inside and outside the ground-truth defect
mask, Milestone S2 (MSE) vs.\ Milestone S5 (L1), across all 119 defective
screw test images. Checkpoints predate the final seed-42 models in
Table~\ref{tab:screw-results}.}
\begin{tabular}{lccc}
\toprule
Model & Defect Error (signal) & Background Error (noise) & SNR \\
\midrule
Milestone S2 (MSE) & 0.018015 & 0.001330 & 13.54 \\
Milestone S5 (L1 loss) & 0.024089 & 0.001447 & 16.65 \\
\bottomrule
\end{tabular}
\label{tab:screw-snr}
\end{table}

\subsubsection{Summary of Screw Exploration Results}

Table~\ref{tab:screw-results} summarises the full screw exploration phase.
All autoencoder rows use seed 42, a 20-epoch training budget, and the 85/15
split described above; the baseline row is the mean of the three-seed check.
Pixel F1 and Image F1 use the same validation-derived thresholds as the
adjacent accuracy columns (99th percentile pixel-level, 95th percentile
image-level). Milestones S3 and S4's values should be read as a single-run
snapshot affected by the diagnosed seed-instability issue in their masking
implementation (unseeded \texttt{random.randint()}, not covered by
\texttt{torch.manual\_seed()}): as Milestone S3's write-up above shows, a
different unseeded rerun of the same code can produce materially different
numbers for this specific pair. Every other row in this table is seed-stable.

\begin{table}[h]
\centering
\caption{Screw-category results across the exploratory autoencoder variants.}
\resizebox{\textwidth}{!}{%
\begin{tabular}{lcccccccc}
\toprule
Model & Pixel AP & Pixel Acc & Pixel F1 & Top-1\% Capture & Pixel AUROC & Img F1 & Img MSE AUROC & Top-K AUROC \\
\midrule
Z-score baseline & 0.0123 & 0.9843 & -- & -- & 0.6887 & -- & -- & -- \\
Milestone S1 & 0.0225 & 0.9885 & 0.0557 & 0.1417 & 0.8723 & 0.000 & 0.8541 & 0.9561 \\
Milestone S2 & 0.0453 & 0.9872 & 0.0900 & 0.2269 & 0.8849 & 0.168 & 0.5431 & 0.7225 \\
Milestone S3 & 0.0117 & 0.9868 & 0.0320 & 0.0852 & 0.7765 & 0.025 & 0.0047 & 0.3210 \\
Milestone S4 & 0.0173 & 0.9826 & 0.0526 & 0.2015 & 0.7401 & 0.053 & 0.0842 & 0.0502 \\
\textbf{Milestone S5 (L1 loss)} & \textbf{0.0581} & 0.9871 & 0.1029 & \textbf{0.2508} & \textbf{0.9075} & 0.126 & 0.8176 & 0.8721 \\
\bottomrule
\end{tabular}%
}
\label{tab:screw-results}
\end{table}

Image F1 of 0.000 for Milestone S1 is a genuine measurement at the
validation-derived operating threshold, not a data error. It means that this
operating point produced no true-positive image detections on the screw test
set. The result should therefore be interpreted together with the
threshold-independent image-level AUROC, which shows that the model can retain
ranking information even when its chosen operating threshold performs poorly.

These exploratory results predate and motivated the production PatchCore and
Keras CAE baseline comparison reported in Section~\ref{sec:evaluation}, where screw achieves a
PatchCore image F1 of 0.796 and pixel AUPIMO of 0.380, and a Keras CAE image
F1 of 0.421 and pixel AUPIMO of 0.005 (Table~\ref{tab:final_results}), figures that already reflect
the lessons of this exploration: screw's rotational variability and greyscale
presentation make it one of the harder categories in the benchmark,
consistent with what we found here.

The gap is largest on image-level F1 specifically: Milestone S5, the best
model in this exploration, reaches an image F1 of only 0.126, well short of
both the PatchCore (0.796) and Keras CAE (0.421) screw baselines above. This
is expected rather than a weakness specific to screw: none of these
exploratory autoencoders use a pretrained backbone, and the same
from-scratch-versus-pretrained-features gap is already visible in this
report's own bottle exploration (Section~\ref{sec:ali_baseline}), where a frozen-ImageNet
5-NN baseline (AUROC 0.988) decisively outperformed every from-scratch
autoencoder variant (best AUROC 0.877) and directly motivated adopting
PatchCore. The screw milestones here should be read the same way: as
exploratory evidence about masking, compression, and loss-function choices
on a from-scratch architecture, not as a competitive alternative to the
production models above.

%
% Figures:
%   figures/screw/fig_v2_v4_residual_snr.png
%   figures/screw/fig_v2_localization_examples.png
%   figures/screw/fig_v3_masking_bug_vs_fix.png
%   figures/screw/figures/masking_example_background_miss.png
%   figures/screw/figures/masking_example_object_hit.png
%
% Cross-references below use the existing report labels rather than
% hardcoded section/table numbers.
%
% Note for the main report:
%   Section 2.2 / 4.2 still need the separate orientation wording update
%   so that screw is treated as a random-rotation object-category exception.
% ============================================================

\subsection{Parallel Screw-Category Exploration}
\label{subsec:screw-exploration}

Alongside the bottle-category exploration described above, we conducted a
parallel iterative development phase on the screw category. Screw presents
two structural challenges not shared by bottle: its images are
single-channel greyscale rather than RGB, and, per the original MVTec AD
benchmark, its parts are photographed at random rotation on the inspection
surface rather than arriving in a consistent orientation, one of the
reasons the benchmark identifies screw as one of the dataset's most
difficult categories. Under the same deterministic 85\%/15\% protocol used
throughout this report, screw's 320 normal training images are divided
into 272 fitting images and 48 validation images. As with the bottle
milestones above, the official test set (160 images: 41 normal, 119
defective) is used to score each milestone as it is built, not reserved
for a single blind check at the end, so the results below are exploratory
and test-informed in exactly the same sense as Milestones 1-4 for bottle.
We report this phase for the same reason we report those milestones: the
dead ends are as informative as the improvements.

This screw exploration was implemented independently in PyTorch, separately
from the Keras implementation of the final CAE described in
Section~\ref{sec:cae}. The two are not architecturally identical and share no code.

\textbf{A note on naming.} The five variants below are referred to as
Milestones S1-S5 in this write-up, to parallel this report's Milestone
1-5 naming for bottle. The figures accompanying Milestones S3-S5 were
generated directly by the exploration notebook and keep that notebook's
own labels (V1, V2, V3, V3b, V4) baked into the image content. The
mapping is: V1 = S1, V2 = S2, V3 = S3, V3b = S4, V4 = S5. 

Supplementary visualisations from these experiments are collected in
Appendix~\ref{sec:screw-exploration-figures}.

\subsubsection{Baselines: Z-Score Model and the Accuracy Paradox in Practice}

Before any learned model, we built a non-learned baseline: a per-pixel mean and
standard deviation computed from the normal training images, with test-time
pixels flagged when they deviate significantly from that per-pixel statistic,
conceptually identical to the Variation Model benchmarked in the original
MVTec AD paper. A dedicated three-seed robustness check (seeds 42, 17, 100)
confirmed this baseline is stable regardless of the train/validation split
seed, with a Pixel AP standard deviation of 0.0003 and a Pixel AUROC standard
deviation of 0.0018; mean Pixel AP across the three seeds is 0.0123 at 98.43\%
accuracy.

To make the accuracy-paradox argument from Section~\ref{sec:imbalance} concrete on this
category rather than only cited, we additionally scored a trivial
``always normal'' classifier against the true defect masks: it achieved
99.75\% pixel accuracy against just 0.0025 AP. The z-score baseline
deliberately trades a little accuracy for a large jump in AP, the opposite
direction accuracy alone would suggest: under the single seed-42 run used
throughout the rest of this section, it reaches 98.43\% accuracy and 0.0126
AP, consistent with the three-seed mean above, which is the clearest
demonstration in our screw experiments of why AP-style metrics are
necessary here.

\subsubsection{Milestone S1: Baseline Convolutional Autoencoder}

Our first screw architecture used a standard encoder-decoder with an implicit
spatial bottleneck (no explicit dense compression layer), trained with MSE
reconstruction loss on normal-only images. This establishes the reference
point every later change is measured against. The result was a Pixel AP of
0.0225, a real but modest signal that confirmed the pipeline worked end to end
without yet being competitive.

\subsubsection{Milestone S2: Explicit 128-Dimensional Bottleneck}

We added a dense linear bottleneck, forcing a much stronger, more deliberate
compression than Milestone S1's implicit spatial one. The result was a Pixel
AP of 0.0453, roughly double Milestone S1's. This was the first strong
confirmation on screw that compression strength is a real, controllable lever
for anomaly separation, not just capacity for cleaner reconstruction, the same
conclusion Section~\ref{sec:ae3} reaches independently for bottle.

\subsubsection{Milestone S3: Masked-Restoration Self-Supervision and the Overlap Diagnostic}

We tested a different strategy for the same underlying goal explored in
Milestone S2, preventing the model from trivially learning to copy its input,
adapted from the masked-restoration technique introduced for bottle in
Milestone 2 (Section~\ref{sec:ae2}): random patches were zeroed and the model
was trained to restore the full image. This failed decisively
on screw. We diagnosed the failure quantitatively rather than assuming it: a
targeted check sampling 500 random $32\times32$ training patches found that
only 190 (38.0\%) touched the screw object at all, meaning 62.0\% of masked
patches during training were pure background. Most of the training signal
taught the model nothing about screw structure. The result was a Pixel AP of
0.0117, the worst model in this phase, below even the non-learned z-score
baseline (0.0123).

Milestone S3 uses Python's built-in random module for masked-patch placement.
It calls \texttt{random.randint()}, which is not controlled by
\texttt{torch.manual\_seed()}. Because \texttt{random.seed()} is never set before this milestone trains, its patch
locations, and therefore its trained weights and final metrics, depend on
whatever random state the kernel happens to start with: across two
independent reruns of identical code, Pixel AP ranged from 0.0146 to 0.0394
and image-level AUROC ranged from 0.000 to 0.563. Milestone S4 (below) uses
the same unseeded function but happens to run after an unrelated
diagnostic cell that calls \texttt{random.seed(0)} as a side effect, which
is why it reproduces exactly across reruns despite sharing the same
underlying gap; that reproducibility is an artefact of cell execution
order, not an intentional seed. The fix is to add
\texttt{random.seed(RANDOM\_SEED)} alongside the existing PyTorch seed.

Because this milestone's patch placement is unseeded, where the mask lands
varies from run to run rather than being fixed by the model or data. Two
additional single-patch examples illustrate this directly: in one run the
masked patch sits almost entirely on background, away from the screw body;
in another run of the identical code, the patch lands closer to the object,
overlapping its edge. This is the same unseeded call landing differently
each time the kernel starts, not two different bugs, which is exactly why
this milestone's metrics vary run to run rather than the placement being
consistently bad or consistently fine.

\subsubsection{Milestone S4: Object-Centred Masking Fix}

We directly addressed the diagnosed Milestone S3 bug by restricting mask
placement to the object's foreground region, using the same
foreground-extraction pipeline described in Section~\ref{sec:foreground}, rather than masking
uniformly at random across the full frame. The bug fix itself is correct and
verifiable independent of any model score. However, fixing the bug did not
translate into a competitive result, but it did help: this model's Pixel AP
(0.0173) is roughly 1.5x Milestone S3's (0.0117) and moves back above the
non-learned z-score baseline (0.0123), though it still falls well short of
the bottleneck-based milestones (S1, S2, S5). We consider this an important
negative-then-partial-recovery result of our screw experiments: it
demonstrates that a correct bug diagnosis and a correct fix do not
automatically make the underlying method competitive.

\subsubsection{Milestone S5: L1 Loss Ablation}

Our final screw variant kept Milestone S2's architecture unchanged and
replaced the MSE reconstruction loss with an L1 (mean absolute error) loss,
motivated by the hypothesis that L1's sharper, less blur-prone gradient
behaviour would produce cleaner separation between normal and defective pixel
errors. This produced the best result of the screw exploration phase: Pixel
AP 0.0581, Pixel AUROC 0.9075, and a top-1\% defect-pixel capture rate of
0.2508.

To rule out the simplest alternative explanation, that L1 was simply
elevating error everywhere including the background rather than genuinely
separating defect regions better, we measured mean reconstruction error
inside and outside the ground-truth defect mask across all 119 defective
screw test images (Table~\ref{tab:screw-snr}). The L1 model's background
error is only marginally higher than the MSE model's, while its
defect-region error is proportionally much higher, giving it a higher
signal-to-noise ratio and ruling out the simple noise-amplification
explanation. We note this analysis was run against an earlier,
pre-final checkpoint pair, so the qualitative finding is suggestive rather
than confirmed against the exact final models in
Table~\ref{tab:screw-results}.

\begin{table}[h]
\centering
\caption{Mean reconstruction error inside and outside the ground-truth defect
mask, Milestone S2 (MSE) vs.\ Milestone S5 (L1), across all 119 defective
screw test images. Checkpoints predate the final seed-42 models in
Table~\ref{tab:screw-results}.}
\begin{tabular}{lccc}
\toprule
Model & Defect Error (signal) & Background Error (noise) & SNR \\
\midrule
Milestone S2 (MSE) & 0.018015 & 0.001330 & 13.54 \\
Milestone S5 (L1 loss) & 0.024089 & 0.001447 & 16.65 \\
\bottomrule
\end{tabular}
\label{tab:screw-snr}
\end{table}

\subsubsection{Summary of Screw Exploration Results}

Table~\ref{tab:screw-results} summarises the full screw exploration phase.
All autoencoder rows use seed 42, a 20-epoch training budget, and the 85/15
split described above; the baseline row is the mean of the three-seed check.
Pixel F1 and Image F1 use the same validation-derived thresholds as the
adjacent accuracy columns (99th percentile pixel-level, 95th percentile
image-level). Milestones S3 and S4's values should be read as a single-run
snapshot affected by the diagnosed seed-instability issue in their masking
implementation (unseeded \texttt{random.randint()}, not covered by
\texttt{torch.manual\_seed()}): as Milestone S3's write-up above shows, a
different unseeded rerun of the same code can produce materially different
numbers for this specific pair. Every other row in this table is seed-stable.

\begin{table}[h]
\centering
\caption{Screw-category results across the exploratory autoencoder variants.}
\resizebox{\textwidth}{!}{%
\begin{tabular}{lcccccccc}
\toprule
Model & Pixel AP & Pixel Acc & Pixel F1 & Top-1\% Capture & Pixel AUROC & Img F1 & Img MSE AUROC & Top-K AUROC \\
\midrule
Z-score baseline & 0.0123 & 0.9843 & -- & -- & 0.6887 & -- & -- & -- \\
Milestone S1 & 0.0225 & 0.9885 & 0.0557 & 0.1417 & 0.8723 & 0.000 & 0.8541 & 0.9561 \\
Milestone S2 & 0.0453 & 0.9872 & 0.0900 & 0.2269 & 0.8849 & 0.168 & 0.5431 & 0.7225 \\
Milestone S3 & 0.0117 & 0.9868 & 0.0320 & 0.0852 & 0.7765 & 0.025 & 0.0047 & 0.3210 \\
Milestone S4 & 0.0173 & 0.9826 & 0.0526 & 0.2015 & 0.7401 & 0.053 & 0.0842 & 0.0502 \\
\textbf{Milestone S5 (L1 loss)} & \textbf{0.0581} & 0.9871 & 0.1029 & \textbf{0.2508} & \textbf{0.9075} & 0.126 & 0.8176 & 0.8721 \\
\bottomrule
\end{tabular}%
}
\label{tab:screw-results}
\end{table}

Image F1 of 0.000 for Milestone S1 is a genuine measurement at the
validation-derived operating threshold, not a data error. It means that this
operating point produced no true-positive image detections on the screw test
set. The result should therefore be interpreted together with the
threshold-independent image-level AUROC, which shows that the model can retain
ranking information even when its chosen operating threshold performs poorly.

These exploratory results predate and motivated the production PatchCore and
Keras CAE baseline comparison reported in Section~\ref{sec:evaluation}, where screw achieves a
PatchCore image F1 of 0.796 and pixel AUPIMO of 0.380, and a Keras CAE image
F1 of 0.421 and pixel AUPIMO of 0.005 (Table~\ref{tab:final_results}), figures that already reflect
the lessons of this exploration: screw's rotational variability and greyscale
presentation make it one of the harder categories in the benchmark,
consistent with what we found here.

The gap is largest on image-level F1 specifically: Milestone S5, the best
model in this exploration, reaches an image F1 of only 0.126, well short of
both the PatchCore (0.796) and Keras CAE (0.421) screw baselines above. This
is expected rather than a weakness specific to screw: none of these
exploratory autoencoders use a pretrained backbone, and the same
from-scratch-versus-pretrained-features gap is already visible in this
report's own bottle exploration (Section~\ref{sec:ali_baseline}), where a frozen-ImageNet
5-NN baseline (AUROC 0.988) decisively outperformed every from-scratch
autoencoder variant (best AUROC 0.877) and directly motivated adopting
PatchCore. The screw milestones here should be read the same way: as
exploratory evidence about masking, compression, and loss-function choices
on a from-scratch architecture, not as a competitive alternative to the
production models above.
